 \documentclass[leqno, a4paper, 12pt]{jsarticle}
\usepackage{amssymb}
\usepackage{amsthm}
\usepackage{amsmath}
\usepackage{comment}
\usepackage{url}
	%可換図式用
		%\usepackage[all]{xy}
	%重くなるので使わいときはコメント化しておく。
%定理環境 thmはあまり使わずpropをなるべく使う。
%主張は基本的に証明の中で使い、そのため番号を付けない。
%注意環境を付けてみた。
\newtheorem{thm}{定理}
\newtheorem{prop}[thm]{命題}
\newtheorem{cor}[thm]{系}
\newtheorem{lem}[thm]{補題}
\newtheorem{df}[thm]{定義}
\newtheorem{exam}[thm]{例}
\newtheorem{fact}[thm]{事実}
\newtheorem*{claim}{主張}
\newtheorem*{rmk}{注意}
\renewcommand{\proofname}{{\bf 証明}}
%矢印たち。
%\toは\rightarrowと同じ。\getsは\leftarrowと同じ。同値は\iff
%上向き矢はuparrow逆はdownarrow
\newcommand{\To}{\Longrightarrow}
\newcommand{\Gets}{\Longleftarrow}
%黒板太字
\newcommand{\nn}{\mathbb{N}}
\newcommand{\zz}{\mathbb{Z}}
\newcommand{\qq}{\mathbb{Q}}
\newcommand{\rr}{\mathbb{R}}
\newcommand{\cc}{\mathbb{C}}
%無限大
\newcommand{\mgn}{\infty}
%ギリシャン
\newcommand{\dpst}{\displaystyle}
\newcommand{\ep}{\varepsilon}
\newcommand{\al}{\alpha}
\newcommand{\be}{\beta}
\newcommand{\de}{\delta}
\newcommand{\la}{\lambda}
\newcommand{\La}{\Lambda}
\newcommand{\ga}{\gamma}
\newcommand{\lil}{\la\in \La}
%商集合
\newcommand{\qt}[2]{#1/\! #2}
%enumerate環境
\renewcommand{\labelenumi}{{\rm (\arabic{enumi})}}
%以降alignじゃなくてenumerate を使え。
%なんか演算
\DeclareMathOperator{\card}{card}
\DeclareMathOperator{\supp}{supp}
%なんか演算２
\DeclareMathOperator{\bd}{Bdry}
\DeclareMathOperator{\cl}{CL}
\DeclareMathOperator{\nai}{Int}
\DeclareMathOperator{\di}{diam}
\DeclareMathOperator{\pr}{pr}

\DeclareMathOperator{\rank}{rank}
%差集合は\setminusで統一する。等号の否定は\neq
%真部分集合は\subsetneqq　偏微分は\partial
%ディスプレイスタイル。\dfracでディスプレイ分数
%空集合は\Oを使う！！！！！！！！！！！！

\newcounter{mycnt}%                       カウンター定義
\newcommand{\outcnt}{%                    コマンド定義
  \the\value{mycnt}%                      カウンター出力
  \ifnum\value{mycnt}>100 (over 100)\fi%  100以上の場合
  \stepcounter{mycnt}%                    mycnt++
}
\setcounter{mycnt}{1}

\title{
多項式に関する一致の定理
}
\author{ナンブキトラ}
\date{\today}
			\begin{document}
\maketitle

この文章では多項式に関する一致の定理を証明する．
つまり実係数多項式$p$の零点全体の集合が内点を持つならば
恒等的に$0$に等しいことを証明する．

まず最初に初等的な行列式の計算を行う．
\begin{thm}[ファンデルモンドの行列式]
$n\in \zz_{\ge 1}$とする．
このとき任意の$x_{1}\dots, x_{n}\in \rr$
に対して
\[
\begin{vmatrix}
1 & x_{1} & x_{1}^{2} & \dots & x_{1}^{n-1}\\
1 & x_{2} & x_{2}^{2} & \dots & x_{2}^{n-1}\\
\vdots & \vdots  &  \vdots & \dots &\vdots\\
1 & x_{n} & x_{n}^{2} & \dots & x_{n}^{n-1}
\end{vmatrix}
=\prod_{i<j}(x_{j}-x_{i})
\]
が成り立つ．
\end{thm}
\begin{proof}
$n$に関する帰納法で証明する．
$n+1$より小さい数では成り立つと仮定する．
このとき第$1$行を各行から引くと
\begin{align*}
&
\begin{vmatrix}
1 & x_{1} & x_{1}^{2} & \dots & x_{1}^{n-1} & x_{1}^{n}\\
1 & x_{2} & x_{2}^{2} & \dots & x_{2}^{n-1} & x_{2}^{n}\\
\vdots & \vdots  &  \vdots & \dots &\vdots & \vdots\\
1 & x_{n} & x_{n}^{2} & \dots & x_{n}^{n-1} & x_{n}^{n}\\
1 & x_{n+1} & x_{n+1}^{2} & \dots & x_{n+1}^{n-1} & x_{n+1}^{n}\\
\end{vmatrix}
\\
&=
\begin{vmatrix}
1 & x_{1} & x_{1}^{2} & \dots & x_{1}^{n-1} & x_{1}^{n}\\
0 & x_{2}-x_{1} & x_{2}^{2}-x_{1}^{2} & \dots & x_{2}^{n-1}-x_{1}^{n-1} & x_{2}^{n}-x_{1}^{n}\\
\vdots & \vdots  &  \vdots & \dots &\vdots & \vdots\\
0 & x_{n}-x_{1} & x_{n}^{2} -x_{1}^{2}& \dots & x_{n}^{n-1}-x_{1}^{n-1} & x_{n}^{n}-x_{1}^{n}\\
0 & x_{n+1} -x_{1}& x_{n+1}^{2}-x_{1}^{2} & \dots & x_{n+1}^{n-1} -x_{1}^{n-1}& x_{n+1}^{n}-x_{1}^{n}\\
\end{vmatrix}
\end{align*}
となる．計算は次のページに続く

さらに行列式の展開の性質から
\begin{align*}
=
\begin{vmatrix}
 x_{2}-x_{1} & x_{2}^{2}-x_{1}^{2} & \dots & x_{2}^{n-1}-x_{1}^{n-1} & x_{2}^{n}-x_{1}^{n}\\
\vdots & \vdots  &  \vdots & \dots &\vdots & \vdots\\
 x_{n}-x_{1} & x_{n}^{2} -x_{1}^{2}& \dots & x_{n}^{n-1}-x_{1}^{n-1} & x_{n}^{n}-x_{1}^{n}\\
 x_{n+1} -x_{1}& x_{n+1}^{2}-x_{1}^{2} & \dots & x_{n+1}^{n-1} -x_{1}^{n-1}& x_{n+1}^{n}-x_{1}^{n}\\
\end{vmatrix}
\end{align*}
となる．
ここで第$k$列の$-x_{1}$倍をを$k+1$列に足すと最後の両列式は
\begin{align*}
&=
\begin{vmatrix}
 x_{2}-x_{1} & x_{2}^{2}-x_{1}^{2} & \dots & x_{2}^{n-1}-x_{1}^{n-1} & x_{2}^{n}-x_{1}^{n}\\
\vdots & \vdots  &  \vdots & \dots &\vdots & \\
 x_{n}-x_{1} & x_{n}^{2} -x_{1}^{2}& \dots & x_{n}^{n-1}-x_{1}^{n-1} & x_{n}^{n}-x_{1}^{n}\\
 x_{n+1} -x_{1}& x_{n+1}^{2}-x_{1}^{2} & \dots & x_{n+1}^{n-1} -x_{1}^{n-1}& x_{n+1}^{n}-x_{1}^{n}\\
\end{vmatrix}
\\
&=
\begin{vmatrix}
 x_{2}-x_{1} & x_{2}^{2}-x_{1}x_{2} & \dots & x_{2}^{n-1}-x_{1}x_{2}^{n-2} & x_{2}^{n}-x_{1}x_{n}^{n-1}\\
\vdots & \vdots  &  \vdots & \dots &\vdots & \\
 x_{n}-x_{1} & x_{n}^{2} -x_{1}x_{n}& \dots & x_{n}^{n-1}-x_{1}x_{n}^{n-2} & x_{n}^{n}-x_{1}x_{n}^{n-1}\\
 x_{n+1} -x_{1}& x_{n+1}^{2}-x_{1}x_{n} & \dots & x_{n+1}^{n-1} -x_{1}x_{n}^{n-2}& x_{n+1}^{n}-x_{1}x_{n}^{n-1}\\
\end{vmatrix}
\\
&=
\begin{vmatrix}
 x_{2}-x_{1} & x_{2}(x_{2}-x_{1}) & \dots & x_{2}^{n-2}(
 x_{2}-x_{1}) & x_{2}^{n-1}(x_{2}-x_{1})\\
\vdots & \vdots  &  \vdots & \dots &\vdots & \\
 x_{n}-x_{1} & x_{n}(x_{n} -x_{1})& \dots & x_{n}^{n-2}(x_{n}-x_{1}) & x_{n}^{n-1}(x_{n}-x_{1})\\
 x_{n+1} -x_{1}& x_{n+1}(x_{n+1}-x_{1}) & \dots & x_{n+1}^{n-2}(x_{n+1} -x_{1})& x_{n+1}(x_{n+1}-x_{1})\\
\end{vmatrix}
\\
&=
\prod_{1<j}(x_{j}-x_{1})\cdot 
\begin{vmatrix}
1 & x_{2} & x_{2}^{2} & \dots & x_{2}^{n-1}\\
1 & x_{3} & x_{3}^{2} & \dots & x_{3}^{n-1}\\
\vdots & \vdots  &  \vdots & \dots &\vdots\\
1 & x_{n+1} & x_{n+1}^{2} & \dots & x_{n+1}^{n-1}
\end{vmatrix}
=\prod_{i<j}(x_{j}-x_{i})
\end{align*}
となるので証明が終わる．
\end{proof}
\begin{proof}[\textbf{別証明}]
与えられた行列式を$p_{n}(x_{1}, \dots x_{n})$
とするとこれは$x_{1}\dots, x_{n}$に関する多項式で
しかも各$x_{i}$に関する最高次数は$n-1$である．
さて，行列式の性質から$x_{i}=x_{j}$とすると
$p_{n}=0$となる．
ゆえに$p_{n}$は$x_{j}-x_{i}$を因数にもつ．
よって多項式$q_{n}=\prod_{i<j}(x_{j}-x_{i})$
によって$p_{n}$を割り切ることができる．ここで
多項式$q_{n}$について各$x_{i}$の最高次数は$n-1$であるから
$p_{n}$は$q_{n}$の定数倍であることがわかる．
つまり$p_{n}=a_{n}q_{n}$となる$a_{n}\in \rr$が存在する．
簡単に$a_{1}=1$がわかり，また
$n+1$の場合に$x_{1}=0$とすると
\begin{align*}
p_{n+1}&=
\begin{vmatrix}
1 & 0 & 0 & \dots & 0 & 0\\
1 & x_{2} & x_{2}^{2} & \dots & x_{2}^{n-1} & x_{2}^{n}\\
\vdots & \vdots  &  \vdots & \dots &\vdots & \vdots\\
1 & x_{n} & x_{n}^{2} & \dots & x_{n}^{n-1} & x_{n}^{n}\\
1 & x_{n+1} & x_{n+1}^{2} & \dots & x_{n+1}^{n-1} & x_{n+1}^{n}\\
\end{vmatrix}
=\prod_{j>1}x_{j}\cdot
\begin{vmatrix}
1 & x_{2} & x_{2}^{2} & \dots & x_{2}^{n-1}  \\
\vdots & \vdots  &  \vdots & \dots &\vdots \\
1 & x_{n} & x_{n}^{2} & \dots & x_{n}^{n-1} \\
1 & x_{n+1} & x_{n+1}^{2} & \dots & x_{n+1}^{n-1} \\
\end{vmatrix}
\\
&=
\left(\prod_{j>1}x_{j}\right)\cdot
p_{n}(x_{2}, \dots x_{n+1})
\end{align*}
となり，また
\[
q_{n+1}(0, x_{2}, \dots x_{n})
=\left(\prod_{j>1}x_{j}\right)\cdot
q_{n}( x_{2}, \dots x_{n})
\]
なので$a_{n+1}=a_{n}$がわかり，結局$a_{n}=1$
がわかるので，定理は成り立つ．
\end{proof}


\begin{thm}[多項式に関する一致の定理]
$p$を$n$変数の実係数多項式とする．
もしも$p$の零点全体の集合$Z$が内点を持つならば
$p$は恒等的に$0$となる．
つまり$p$が定数$0$に等しくないならば
$p(x)\neq 0$となる$x\in \rr^{n}$全体の集合は$\rr^{n}$
のなかで稠密になる．
\end{thm}



\begin{proof}[\textbf{証明その\outcnt}]
$n$に関する帰納法で証明する．$n+1$より小さい数では
命題が成り立つと仮定する．
$p(x_{1}, \dots, x_{n}, x_{n+1})$
についてこれは多項式なので，
$x_{n+1}$に関する最高次数を$m$として
$x=(x_{1}, \dots, x_{n})$として以下のように表すことができる．
\[
p(x_{1}, \dots, x_{n}, x_{n+1})=
\sum_{i=0}^{m}q_{i}(x)x_{n+1}^{i}
\]
ここで各$q_{i}$は$x$の多項式である．
今$p$の零点全体を$Z$と書くと，
$Z$は内点を持つので開集合$V\subset \rr^{n}$
と開集合$U\subset \rr$が存在して
$V\times U\subset Z$となる．
$U$から相異なる点$t_{1}, \dots, t_{m+1}\in U$を取ると
$k\in \{1, \dots, m+1\}$と$x\in V$について
\[
\sum_{i=0}^{m}q_{i}(x)t_{k}^{i}=
p(x, t)=0
\]
となる．これは次のように書くことができる．
\[
\begin{pmatrix}
1 & t_{1} & t_{1}^{2} & \dots & t_{1}^{m}\\
1 & t_{2} & t_{2}^{2} & \dots & t_{2}^{m}\\
\vdots & \vdots  &  \vdots & \dots &\vdots\\
1 & t_{n} & t_{m}^{2} & \dots & t_{m+1}^{m}
\end{pmatrix}
\begin{pmatrix}
q_{0}\\
q_{1}\\
\vdots\\
q_{m}
\end{pmatrix}
=
\begin{pmatrix}
0\\
0\\
\vdots\\
0
\end{pmatrix}
\]
先の定理から左辺に出てくる行列は可逆であることがわかるので，結局任意の$x\in U$と$i\in \{1, \dots, m\}$について
$q_{i}=0$がわかる．
よって各$q_{i}$の零点全体は内点を持つことがわかったので
帰納法の仮定より，
各$q_{i}$は恒等的に$0$に等しい．
よって$p$も恒等的に$0$に等しい．
\end{proof}


\begin{proof}[\textbf{証明その\outcnt}]
上の証明と同じく帰納法で証明する．
同じ記号を用いる．
$Z$の元を$\rr^{n+1}\setminus Z$の元で近似できれば良い．
$(x, t)\in Z\subset \rr^{n}\times \rr$とする．
この時
\[
\sum_{i=0}^{m}q_{i}(x)t^{i}=0
\]
となる．

Case1 : 
ある$k\in \{0, \dots, m\}$について$q_{k}(x)\neq 0$
の場合．
このときは
$p(x, t)$を$t$の多項式として見れば一変数なので
一変数の場合から$t$に十分近い$t'\in \rr$であって
$p(x, t')\neq 0$となるものが存在する．
よって$(x, t')$が求める$\rr^{n+1}\setminus Z$の元である．


Case2: 全ての$i\in \{1, \dots, m\}$について
$x$が$q_{i}$の零点だった場合．
このときには帰納法の仮定より$x$に十分近い各$q_{i}$の
非零点になっているような$x'\in \rr^{n}$が存在する．
$p(x', t)\neq 0$の場合は$(x', t)\in \rr^{n+1}\setminus Z$
が求めるものになっている．
$p(x', t)=0$の場合にはCase1に還元される．
\end{proof}



\begin{proof}[\textbf{証明その\outcnt}]
$n=1$の時には$n$次多項式は高々$n$個の零点を持たないことから定理が成り立つことがわかる．
次に一般の$n$について考える．
$Z$の内点を$a$とする．
このとき十分小さい$r\in \rr$が存在して
開円盤$B=
\{\, x\in \rr^{n}
\mid 
\|x-a\|_{2}<r
\, \}$
は$Z$の部分集合となる．
ここで
$u\in B$
を任意にとり
$q:\rr\to \rr$を$q(t)=p(a+u\cdot t)$
と定義する．
すると$q$は$t$に関する多項式であり，
$p$が$B$上で$0$になることから$q$は
$t\in (-1, 1)$の時$q(t)=0$となる．
よって$1$変数の場合から$q$は恒等的に$0$
になる．$u\in B$は任意だったので
$p$は全空間で恒等的に$0$になる．
\end{proof}


\begin{proof}[\textbf{証明その\outcnt}]
まず最初に$x\in \rr^{n}$, $y\in \rr$として
多項式$P(x, y)$は$y$の一次式で割り算できることに注意しよう．
さて，$p(x, y)$の零点が内点を持つとする.
内点$(a, b)$について$p(x, b)$
は$n$変数として見たときに内点をもつよって帰納法の，仮定より$p(x, b)$は恒等的に$0$になる．
つまり$y=b$は$p(x, y)$の零点になるので
$p(x, y)$は$(y-b)$を因数に持つ．
$p(x, y)$の次数を$m$として
相異なる点$b_{1}, \dots, b_{m+1}$を
任意の$x$について$p(x, b_{i})=0$
となるようにとる
と
$n$変数多項式$Q$が存在して
$p(x, y)=(y-b_{1})\dots (y-b_{m})Q(x)$
となる．ここで$y=b_{m+1}$とすると
$0=(b_{m+1}-b_{1})\dots (b_{m+1}-b_{m})Q(x)$
なので$Q$は恒等的に$0$になる．
つまり$p$も恒等的に$0$
になる．
\end{proof}

\begin{proof}[\textbf{証明その\outcnt}]
多項式は解析関数だから解析関数の一致の定理から従う．
解析関数に関する一致の定理は
\cite{self}や\cite{hako}
を参照のこと．
\end{proof}













	
\begin{thebibliography}{9}
\bibitem{1}
StackExchangeの質問
``The complement of a zero-set of a polynomial is dense'' 
とその回答
\begin{verbatim}
https://math.stackexchange.com/questions/2025978/
the-complement-of-a-zero-set-of-a-polynomial-is-dense
\end{verbatim}

\bibitem{2}
StackExchangeの質問
``Why if $p(x_{1}, \dots, x_{n})$ is 
polynomial on $\rr^{n}$, 
then $p(x)\neq 0$ is satisfied open dens set?
'' 
とその回答

\begin{verbatim}
https://math.stackexchange.com/questions/1731267/
why-if-p-x-1-dots-x-n-is-polynomial-on-mathbbrn-then-p-x-neq-0/1731288
\end{verbatim}

\bibitem{self}
はてなブログ「電波通信」の記事「実解析関数に関する逆関数定理」
\url{
https://concious4410.hatenablog.com/entry/2021/02/28/130208}
\bibitem{hako}
箱さんのwebページの「解析関数のノート」
\url{https://o-ccah.github.io/docs/analytic-function.html}

\end{thebibliography}




			\end{document}



\begin{comment}
[集合のやつは$\mid$を使う。]
[真なる包含関係]
\subsetneqq
[可換図式]
\[\xymatrix{
&   & (2) \ar@{=>}[rd]&  &  &  & (5) \ar@{=>}[rd]&\\ 
& (1)\ar@{=>}[ru] \ar@{=>}[rd]&  &(4)\ar@{=>}[r]&(7)& (1)\ar@{=>}[ru] \ar@{=>}[rd]& & (7)&\\ 
&   & (3) \ar@{=>}[ur]&  &  &  &(6) \ar@{=>}[ur]&
}\]

[\bigin \endを使わない方法]

{\prop[aa] aaa}
で
\begin{prop}[aa]
aaa
\end{prop}
と同じ\df \thm \proof でも同様

[直和]
$\amalg$

[同値関係でよく使う波線]
\sim

[引数付きの新命令]
\newcommand{\prn}[1]{\|#1\|_p}

[番号のリセット]

\setcounter{equation}{0}


[字体の変更]

{\rm hogehoge}と使う。

\rm ローマン
\it イタリック
\sl 斜体

[supp]

\newcommand{\supp}{\text{{\rm supp}}}

[中央揃えの改行]

	\begin{eqnarray*}
		hoge1&=&hogehoge
		hoge2&=&hogehofe

	\end{eqnarray*}

&&の部分で揃えられる。


[\tag]
\begin{equation}
f(x) = ax^2 + bx + c \tag{1.2.3}
\end{equation}



[場合分け]

	\begin{cases}
		hoge1 & \text{状況１}\\
		hoge2 & \text{状況２}\\
		・
		・
		・
	\end{cases}



\end{comment}





